% Exact live proof prepared from the production-verified staged proof.
% Editorial dependency names have been replaced by live cross-references.
% No conventional theorem name, novelty claim, or priority claim is attached.

\begin{proof}
The assertion is local for the \(G\)-topology on \(X\).  Fix a strict
affinoid domain \(U=\mathcal M(C)\) in \(X\), and put
\[
                V=f^{-1}(U)=\mathcal M(D),\qquad u=f^*g|_V.
\]
The space \(V\) is strict affinoid because \(f\) is finite.

We first record the point-set consequence of the hypotheses.  The map
\(V\to U\) is bijective.  Indeed, finiteness and surjectivity give at least
one lift of every support prime, while radiciality gives a unique lift and a
purely inseparable extension of residue fields.  An absolute value on a
field has a unique extension to a purely inseparable field extension: if
\(a^{p^r}\) lies in the smaller field, then its value must be
\[
                       |a|=|a^{p^r}|^{1/p^r}.
\]
Thus every Berkovich point has exactly one lift.  The resulting continuous
bijection \(V\to U\) is a homeomorphism because \(V\) is compact and \(U\)
is Hausdorff.

By \cite[Definition~5.1 and \S5.2(3),(4)]{GJR26}, \(u\) has a
\(G\)-cover by strict affinoid domains on which it is a finite real linear
combination of logarithms of analytic units.  Quasicompactness of the
strict affinoid \(V\) gives a finite subcover.  By the
Gerritzen--Grauert theorem \cite[Theorem~1.1]{Temkin2005GG}, each member
of this finite cover has a finite cover by rational domains of \(V\).  We
therefore obtain a finite rational \(G\)-cover \(\{W_\alpha\}_\alpha\) of
\(V\) such that
\begin{equation}
 \label{eq:p061-frobenius-upstairs-rpl}
       u|_{W_\alpha}
       =\sum_{j=1}^{m_\alpha}\lambda_{\alpha j}
                      \log|t_{\alpha j}|,
 \qquad
       \lambda_{\alpha j}\in\mathbf R,\quad
       t_{\alpha j}\in\mathcal O(W_\alpha)^\times .
\end{equation}

Fix one \(W=W_\alpha\).  Choose a rational presentation
\[
 W=\bigl\{y\in V:\ |b_i(y)|\leq |b_0(y)|,
                  \ 1\leq i\leq s\bigr\},
\]
where \(b_0,\ldots,b_s\in D\) generate the unit ideal.  Put
\(a_i=b_i^q\in C\).  Choose \(d_0,\ldots,d_s\in D\) with
\(1=\sum_i d_i b_i\).  Since \(q\) is a power of \(p\), Frobenius gives
\[
             1=\left(\sum_i d_i b_i\right)^q
              =\sum_i d_i^q a_i.
\]
Every \(d_i^q\) belongs to \(C\), and \(C\hookrightarrow D\) is injective,
so this is an identity in \(C\).  In particular,
\(a_0,\ldots,a_s\) generate the unit ideal.  Hence
\[
 U_W=\bigl\{x\in U:\ |a_i(x)|\leq |a_0(x)|,
                  \ 1\leq i\leq s\bigr\}
\]
is a rational domain in \(U\).

For the unique \(y\in V\) above \(x\in U\), one has
\(|a_i(x)|=|b_i(y)|^q\).  Consequently
\(f^{-1}(U_W)=W\) as subsets.  The two rational-domain structures agree by
the universal property of rational localization.  The finitely many
domains \(U_{W_\alpha}\) cover \(U\), because the \(W_\alpha\) cover \(V\)
and \(V\to U\) is bijective.  A finite cover by analytic domains is a
\(G\)-cover, so these are legitimate \(G\)-local charts on \(U\).

Write \(F=\mathcal O(U_W)\) and \(E=\mathcal O(W)\).  The hypothesis of the
proposition applied to \(U_W\) gives \(E^q\subseteq F\).  If
\(t\in E^\times\), then
\[
                         T=t^q\in F,
              \qquad S=(t^{-1})^q\in F.
\]
The map \(F\to E\) is injective: the restricted map
\(W\to U_W\) is surjective, so its kernel is nilpotent, while \(F\) is
reduced because \(U_W\) is an analytic domain in the reduced space \(X\).
Since \(TS=1\) in \(E\), it follows that \(TS=1\) in \(F\).  Thus \(T\) is
an analytic unit on \(U_W\), not merely a function with nonzero pointwise
values.

Apply this to every \(t_{\alpha j}\) in
\eqref{eq:p061-frobenius-upstairs-rpl}, writing
\(T_{\alpha j}=t_{\alpha j}^q\).  On \(W_\alpha\),
\[
 \sum_j\lambda_{\alpha j}\log|t_{\alpha j}|
 =f^*\!\left(
       \sum_j\frac{\lambda_{\alpha j}}q\log|T_{\alpha j}|
            \right).
\]
The left side is \(f^*g\).  Surjectivity
\(W_\alpha\to U_{W_\alpha}\) therefore yields
\[
       g|_{U_{W_\alpha}}
       =\sum_j\frac{\lambda_{\alpha j}}q
                         \log|T_{\alpha j}|.
\]
Real coefficients are divisible by \(q\).  These expressions on the
rational \(G\)-cover prove that \(g|_U\) is \(\mathbf R\)-PL.  Since strict
affinoid domains \(U\) form a \(G\)-cover of \(X\),
\cite[\S5.2(4)]{GJR26} finishes the proof.
\end{proof}

\begin{theorem}
\label{thm:p061-simplicial-toric-ph-rpl}
Let \(K\) be a complete non-Archimedean field with nontrivial absolute
value, let \(N\) be a lattice, and let
\(\sigma\subset N_{\mathbf R}\) be a strongly convex simplicial rational
polyhedral cone, not necessarily full dimensional.  Let
\(X_{\sigma,N}\) be the associated split affine toric \(K\)-variety, and
let \(U\) be a good strict analytic domain in
\(X_{\sigma,N}^{\mathrm{an}}\).  Then every classically pluriharmonic
function
\[
                              h:U\longrightarrow\mathbf R
\]
is \(\mathbf R\)-PL.  Moreover, \(h\) is harmonic in the sense of
Gubler--Jell--Rabinoff.
\end{theorem}

\begin{proof}
Put
\[
              N_\sigma=N\cap\operatorname{span}_{\mathbf R}(\sigma),
\]
and write \(r=\dim\sigma\).  The subgroup \(N_\sigma\) is saturated in
\(N\): if a nonzero integer multiple of an element of \(N\) lies in the
displayed real span, then the element itself lies in that span.  Thus
\(N/N_\sigma\) is free.  Choose a lattice complement \(N_{\mathrm c}\),
and let \(v_1,\ldots,v_r\) be the primitive ray generators of
\(\sigma\).  Simpliciality says that these vectors are linearly independent
and span \(N_{\sigma,\mathbf R}\).  Therefore
\[
        N'=\mathbf Z v_1\oplus\cdots\oplus\mathbf Z v_r
                      \oplus N_{\mathrm c}\subseteq N
\]
has finite index.  Relative to \(N'\), the cone is smooth, including when
\(r<\operatorname{rank}N\), and
\begin{equation}
 \label{eq:p061-simplicial-smooth-cover}
       X_{\sigma,N'}\simeq
       \mathbf A^r_K\times\mathbf G_{m,K}^{\operatorname{rank}N-r}.
\end{equation}

Let \(M=\operatorname{Hom}(N,\mathbf Z)\) and
\(M'=\operatorname{Hom}(N',\mathbf Z)\).  Then \(M\subseteq M'\) and
\(H=M'/M\) is finite.  Write
\[
 A_K=K[\sigma^\vee\cap M],
 \qquad B_K=K[\sigma^\vee\cap M'].
\]
The inclusion \(A_K\subseteq B_K\) is finite.  Indeed, for every
\(m'\in\sigma^\vee\cap M'\), some positive multiple \(d m'\) lies in
\(M\), and \(\chi^{m'}\) satisfies
\[
                         T^d-\chi^{d m'}=0
\]
over \(A_K\).  The algebra \(B_K\) is of finite type by Gordan's lemma, so
finite type plus integrality gives finiteness.  It is surjective on spectra
by lying over.  The \(H\)-grading of \(B_K\) has degree-zero subring exactly
\(A_K\), hence the finite toric morphism
\begin{equation}
 \label{eq:p061-simplicial-full-cover}
             X_{\sigma,N'}\longrightarrow X_{\sigma,N}
\end{equation}
is the quotient by the finite diagonalizable group \(D(H)\).  No flatness
claim is made or used.

We now split the quotient over a finite extension of \(K\).  If
\(\operatorname{char}K=0\), put \(M_1=M\) and let \(n\) be the exponent
of \(H\).  If \(\operatorname{char}K=p>0\), write canonically
\[
                              H=H_p\oplus H_{p'},
\]
and let \(M_1\subseteq M'\) be the inverse image of \(H_p\).  Then
\[
                       M_1/M=H_p,
             \qquad M'/M_1=H_{p'},
\]
and we let \(n\) be the exponent of \(H_{p'}\).  In either case, choose a
finite Galois extension \(L/K\) containing all \(n\)-th roots of unity;
take \(L=K\) if \(n=1\).  This extension exists because \(T^n-1\) is
separable: in characteristic zero this is automatic, and in characteristic
\(p\) the integer \(n\) is prime to \(p\).  Notice that this uses the
characteristic of \(K\), not its residue characteristic.  A finite extension
of the complete field \(K\) is complete.

Over \(L\), the factorization of \eqref{eq:p061-simplicial-full-cover} is
\begin{equation}
 \label{eq:p061-simplicial-factorization}
 X_{\sigma,N',L}\longrightarrow X_{\sigma,N_1,L}
                         \longrightarrow X_{\sigma,N,L},
\end{equation}
where \(N'\subseteq N_1\subseteq N\) is dual to
\(M\subseteq M_1\subseteq M'\).  The first arrow is the quotient by the
finite constant group obtained by splitting \(D(M'/M_1)\) over \(L\).
In characteristic zero the second arrow is the identity.

Suppose now that \(\operatorname{char}K=p>0\).  Choose
\(q=p^e\) annihilating \(H_p\).  Then \(qM_1\subseteq M\).  If
\[
 A=L[\sigma^\vee\cap M],
 \qquad B=L[\sigma^\vee\cap M_1],
\]
then \(B^q\subseteq A\): for
\(b=\sum_m c_m\chi^m\), Frobenius gives
\[
                         b^q=\sum_m c_m^q\chi^{qm}\in A.
\]
The finite map \(\operatorname{Spec}B\to\operatorname{Spec}A\) is
radicial.  Indeed, if two primes of \(B\) lie above the same prime of
\(A\), then \(b\) in the first prime implies \(b^q\) lies in the base prime
and hence \(b\) lies in the second; symmetry gives equality.  The same
\(q\)-power calculation makes every residue-field extension purely
inseparable.  Finiteness and lying over give surjectivity.

We also record the analytic base-change property needed to apply
\Cref{prop:p061-frobenius-rpl-descent}.  Let \(Z=\mathcal M(C)\) be any
strict affinoid domain in \(X_{\sigma,N,L}^{\mathrm{an}}\), and let
\(Z_1=\mathcal M(D)\) be its inverse image under the second arrow of
\eqref{eq:p061-simplicial-factorization}.  Compatibility of finite
analytification with affinoid base change
\cite[Fact~4.1.16 and Definition~4.2.5]{Temkin2015Berkovich} gives
\begin{equation}
 \label{eq:p061-simplicial-affinoid-base-change}
                              D=C\otimes_A B.
\end{equation}
The ordinary tensor product is already complete: \(B\) is a finite
\(A\)-module, so the right side is a finite \(C\)-module, and a finite
presentation realizes it as a quotient of a finite free \(C\)-module by a
closed submodule.  Every \(z\in D\) is a finite sum
\(z=\sum_i c_i\otimes b_i\), whence
\[
                     z^q=\sum_i c_i^q\otimes b_i^q\in C.
\]
Thus \(D^q\subseteq C\) on every required affinoid domain, and the same
calculation applies after every rational localization.  The algebra \(C\)
is reduced because the split toric variety is geometrically integral and
\(Z\) is an analytic domain.  Surjectivity of the finite radicial map then
makes \(C\to D\) injective.  All hypotheses of
\Cref{prop:p061-frobenius-rpl-descent} are satisfied.

Let \(h:U\to\mathbf R\) be classically pluriharmonic.  By
\cite[Proposition~7.2(6),(7)]{GR25Classical}, its scalar pullback \(h_L\)
to \(U_L\) is classically pluriharmonic.  It is enough to prove that \(h_L\)
is \(\mathbf R\)-PL on every strict affinoid domain \(Z\) in \(U_L\).  Let
\(Z_1\) be the inverse image of \(Z\) under the second arrow of
\eqref{eq:p061-simplicial-factorization}, and let \(V\) be the inverse image
of \(Z_1\) under the first arrow.  Both are strict affinoid because the
arrows are finite.  The space \(V\) is an analytic domain in the smooth
space in \eqref{eq:p061-simplicial-smooth-cover}, hence it is good, strict,
and quasi-smooth.  The same cited proposition shows that the pullback of
\(h_L\) to \(V\) remains classically pluriharmonic.
\Cref{thm:p061-quasismooth-ph-rpl} makes that pullback \(\mathbf R\)-PL.

The target \(Z_1\) is the analytic categorical quotient of \(V\) by the
finite constant group \(G=D(M'/M_1)_L\).  For completeness, this remains
true after restriction to \(Z_1\): if \(C_1=\mathcal O(Z_1)\) and
\(E=\mathcal O(V)\), then the Reynolds idempotent
\[
                         \frac1{|G|}\sum_{\gamma\in G}\gamma
\]
is a bounded projector on \(E\) with image \(E^G=C_1\).  Here
\(|G|\neq0\) in \(L\); it need not be a unit in the valuation ring.  The
projector commutes with the finite affinoid base change.
\Cref{prop:p061-finite-group-rpl-descent-candidate} therefore descends
\(\mathbf R\)-PL regularity from \(V\) to \(Z_1\).

In characteristic \(p>0\),
\Cref{prop:p061-frobenius-rpl-descent} and
\eqref{eq:p061-simplicial-affinoid-base-change} descend it further from
\(Z_1\) to \(Z\).  In characteristic zero the second arrow is the identity.
Thus \(h_L\) is \(\mathbf R\)-PL on a \(G\)-cover of \(U_L\), and hence on
\(U_L\) by \cite[\S5.2(4)]{GJR26}.
\cite[Corollary~5.12]{GJR26} now descends
\(\mathbf R\)-piecewise linearity from \(L\) to \(K\), giving that \(h\) is
\(\mathbf R\)-PL on \(U\).

Finally, \(h\) and \(-h\) are classically psh by the definition of
classical pluriharmonicity.  By \cite[Theorem~4.3]{GR25Classical}, classical
psh and semipositivity agree for \(\mathbf R\)-PL functions.  Both signs of
\(h\) are therefore semipositive, which is precisely
Gubler--Jell--Rabinoff harmonicity by \cite[Definition~7.2]{GJR26}.
\end{proof}
